\chapter{Higher-order functionals}
\label{chap:higher-order-functional}

\begin{quote}
Spec dir:
\passthrough{\lstinline!specs/2026-05-03-phase-11-higher-order-functional/!}.
Version: \passthrough{\lstinline!0.11.0!}.
\end{quote}

\subsection{What this gives you}\label{phase11-what}

A fourth callable arity for \passthrough{\lstinline!gridcalc::func::evaluate!}
that exposes the contracted-scalar biharmonic
\(\nabla^4\psi\) in the integrand. Together with the three arities from
Phase 4, the supported densities are

\begin{lstlisting}[language={C++}]
double f(double psi);                                                            // arity 1
double f(double psi, const Vec3d& grad);                                         // arity 2
double f(double psi, const Vec3d& grad, double lap);                             // arity 3
double f(double psi, const Vec3d& grad, double lap, double biharm);              // arity 4 (new)
\end{lstlisting}

\passthrough{\lstinline!evaluate!} picks the longest matching arity at
compile time and materializes only the derivatives that arity needs.
At arity 4 it materializes \passthrough{\lstinline!gradient(psi)!},
\passthrough{\lstinline!laplacian(psi)!}, and
\passthrough{\lstinline!biharmonic(psi)!} (Phase 8's fused six-term
direct stencil), then forms the integrand pointwise and reduces with
\passthrough{\lstinline!func::integrate!}. The new fourth argument is a
plain \passthrough{\lstinline!double!} per grid point --- it is the
return type of \passthrough{\lstinline!diff::biharmonic!}.

\textbf{Out of scope at Phase 11.} Rank-3 and rank-4 tensor-valued
derivative arguments (the full \(\nabla^3\psi\), \(\nabla^4\psi\)
objects) are deferred to Phase 13/14, where the tensor type and the
contraction expression-template machinery will land together. If your
density only needs the contracted scalars (Laplacian and biharmonic),
the Phase 11 surface is enough.

\subsection{API}\label{phase11-api}

The signature of \passthrough{\lstinline!evaluate!} is unchanged --- the
new arity is a fourth branch in the existing dispatch ladder. Tag
forwarding to \passthrough{\lstinline!func::integrate!} works as before:

\begin{lstlisting}[language={C++}]
namespace gridcalc::func {

template <class F, class Tag = Pairwise>
double evaluate(const core::Field<double>& psi, F&& f, Tag tag = {});

}  // namespace gridcalc::func
\end{lstlisting}

The compile-time error message for an unsupported callable enumerates
all four arities.

\subsection{Worked example: phase-field-crystal free
energy}\label{phase11-pfc}

The phase-field-crystal (PFC) functional is the canonical use case for
this phase. It models an ordered solid as the minimum of a free
energy whose derivative kernel \((q_0^2 + \nabla^2)^2\) selects a
preferred wavelength near \(q_0\):

\[
F[\psi] \;=\; \int \Big[\tfrac{1}{2}\,\psi\,\big[(q_0^2 + \nabla^2)^2 - \epsilon\big]\,\psi \;+\; \tfrac{1}{4}\,\psi^4\Big]\,\dd V.
\]

Expanding \((q_0^2 + \nabla^2)^2 = q_0^4 + 2q_0^2\nabla^2 + \nabla^4\)
gives the four-term form the integrand lambda implements:

\[
F[\psi] \;=\; \int \Big[\tfrac{1}{2}(q_0^4 - \epsilon)\,\psi^2 \;+\; q_0^2\,\psi\,\nabla^2\psi \;+\; \tfrac{1}{2}\,\psi\,\nabla^4\psi \;+\; \tfrac{1}{4}\,\psi^4\Big]\,\dd V.
\]

For the manufactured solution \(\psi(\mathbf{x}) = \sin x\,\sin y\,\sin z\)
on the periodic box \([0, 2\pi)^3\), the derivative chain is
\(\nabla^2\psi = -3\psi\), \(\nabla^4\psi = 9\psi\), and the relevant
integrals are
\(\int\psi^2 = \pi^3\), \(\int\psi^4 = \tfrac{27\pi^3}{64}\).
With \(q_0 = 1\) and \(\epsilon = 0.1\) the analytical functional value is

\[
F_{\text{ref}} \;=\; \Big[\tfrac{1}{2}(1 - 0.1) + (-3) + \tfrac{1}{2}\cdot 9 + \tfrac{27}{256}\Big]\,\pi^3 \;=\; 2.05546875\,\pi^3 \;\approx\; 63.7383.
\]

Calling \passthrough{\lstinline!func::evaluate!} on the discrete \(\psi\):

\begin{lstlisting}[language={C++}]
#include <cmath>
#include <gridcalc/core/EigenAliases.hpp>
#include <gridcalc/core/Field.hpp>
#include <gridcalc/core/Grid.hpp>
#include <gridcalc/func/Functional.hpp>

using gridcalc::core::Field;
using gridcalc::core::Grid;
using gridcalc::core::Vec3d;
using gridcalc::func::evaluate;

constexpr double kPi = 3.14159265358979323846;

const int N = 64;
const double L = 2.0 * kPi;
const double h = L / N;
const double q0  = 1.0;
const double eps = 0.1;

Grid grid(N, N, N, Vec3d(h, h, h));
Field<double> psi(grid, [](double x, double y, double z) {
  return std::sin(x) * std::sin(y) * std::sin(z);
});

auto pfc_density = [q0, eps](double p, const Vec3d&, double lp, double bp) {
  const double q0_sq = q0 * q0;
  const double q0_4  = q0_sq * q0_sq;
  return 0.5 * (q0_4 - eps) * p * p
       + q0_sq * p * lp
       + 0.5   * p * bp
       + 0.25  * p * p * p * p;
};

const double F = evaluate(psi, pfc_density);  // arity 4 selected
\end{lstlisting}

The lambda's \passthrough{\lstinline!const Vec3d&!} parameter is unused
inside the body --- it appears only to select the 4-arg dispatch branch.
At \passthrough{\lstinline!N=64!} the recovered \(F\) matches
\(F_{\text{ref}}\) to better than \(2\%\) relative; the convergence
sweep across \(N \in \{16, 32, 64\}\) shows the
\(\mathcal{O}(h^2)\) rate inherited from the default-order biharmonic
stencil.

\subsection{Verifying with the spectral path}\label{phase11-spectral}

Phase 9 ships
\passthrough{\lstinline!gridcalc::spectral::biharmonic!} as the FFT-based
verification reference (exact for band-limited inputs). For the worked
example above, the same PFC integrand assembled from the spectral
\passthrough{\lstinline!laplacian!} and
\passthrough{\lstinline!biharmonic!} reduces to a value matching
\(F_{\text{ref}}\) to machine precision:

\begin{lstlisting}[language={C++}]
#include <gridcalc/spectral/Derivatives.hpp>
#include <gridcalc/func/Integrate.hpp>

using gridcalc::func::integrate;

const auto lap    = gridcalc::spectral::laplacian(psi);
const auto biharm = gridcalc::spectral::biharmonic(psi);

Field<double> integrand(grid);
for (int k = 0; k < N; ++k)
  for (int j = 0; j < N; ++j)
    for (int i = 0; i < N; ++i)
      integrand(i, j, k) = pfc_density(psi(i, j, k), {}, lap(i, j, k), biharm(i, j, k));

const double F_spectral = integrate(integrand);  // |F_spectral - F_ref| / F_ref < 1e-12
\end{lstlisting}

This is the cross-check the Phase 11 test fixture
(\href{../../../test/pfc_functional_test.cpp}{\passthrough{\lstinline!test/pfc\_functional\_test.cpp!}})
runs as
\passthrough{\lstinline!PfcFunctionalTest.SpectralCrossCheck!} alongside
the FD-path acceptance test
\passthrough{\lstinline!PfcFunctionalTest.HandComputedAtN64!}. The
spectral path requires
\passthrough{\lstinline!-DGRIDCALC\_USE\_FFT=ON!} (the default); the
test \passthrough{\lstinline!\#ifdef!}-skips otherwise.

\subsection{What this does \emph{not} do}\label{phase11-not-do}

\begin{itemize}
\tightlist
\item
  \textbf{Tensor-valued callable arguments.} Rank-3 / rank-4 derivative
  tensors are deferred to Phase 13/14.
\item
  \textbf{The 3rd-derivative slot.} \(\nabla^3\psi\) has no useful
  scalar contraction; the natural rank-1 form
  \(\nabla(\nabla^2\psi)\) is not exposed at Phase 11. PFC does not
  exercise it.
\item
  \textbf{Contraction expression templates.} Deferred to Phase 14
  alongside the tensor-valued callable surface that motivates them.
\item
  \textbf{Variational derivative \(\delta F / \delta \psi\) and PFC
  dynamics.} Phase 11 evaluates the functional value only.
\item
  \textbf{Vector- or tensor-field functionals
  \(F[\mathbf{V}] = \int f(\mathbf{V}, \nabla\mathbf{V}, \dots)\,\dd V\).}
  Phase 14.
\end{itemize}

For the leading-order discrete-error scaling of the biharmonic term, the
spectral cross-check argument, and the rationale for keeping
contraction expression templates out of this phase, see the
``Higher-order functionals'' chapter of the Developer Note.
